\begin{PSbox}[unbreakable]{obj}{Learning Objectives}

This capstone module integrates ALL previous modules into real Electrical Engineering applications. Students will be able to:

\begin{enumerate}
\def\labelenumi{\arabic{enumi}.}
\tightlist
\item
  Apply probabilistic modeling to communication systems
\item
  Analyze signal processing problems statistically
\item
  Model electronic component variations
\item
  Assess power system reliability
\item
  Handle measurement uncertainty in control systems
\item
  Connect probability/statistics to machine learning foundations
\end{enumerate}

\end{PSbox}

\PSsection{16.1}{Communications}

\PSsubsection{16.1.1 BER Estimation and Confidence Intervals}

\textbf{Problem:} Estimate the Bit Error Rate of a communication link.

\textbf{Model:} Each bit is a Bernoulli trial with error probability $p$ (the BER).

\begin{itemize}
\tightlist
\item
  $n$ bits transmitted, $k$ errors observed
\item
  $\hat{p} = \frac{k}{n}$ (MLE of BER)
\item
  95\% CI: $\hat{p} \pm 1.96\sqrt{\frac{\hat{p}(1-\hat{p})}{n}}$
\end{itemize}

\textbf{MATLAB: BER Simulation and CI}

\tcbset{PScodemode/.style={unbreakable}}

\begin{Shaded}
\begin{Highlighting}[]
\CommentTok{\%\% BER Estimation with Confidence Interval}
\VariableTok{EbN0\_dB} \OperatorTok{=} \FloatTok{8}\OperatorTok{;} \VariableTok{EbN0} \OperatorTok{=} \FloatTok{10}\OperatorTok{\^{}}\NormalTok{(}\VariableTok{EbN0\_dB}\OperatorTok{/}\FloatTok{10}\NormalTok{)}\OperatorTok{;}
\VariableTok{N\_bits} \OperatorTok{=} \FloatTok{1e6}\OperatorTok{;}

\CommentTok{\% BPSK transmission}
\VariableTok{bits} \OperatorTok{=} \FloatTok{2}\OperatorTok{*}\VariableTok{randi}\NormalTok{([}\FloatTok{0}\OperatorTok{,}\FloatTok{1}\NormalTok{]}\OperatorTok{,} \FloatTok{1}\OperatorTok{,} \VariableTok{N\_bits}\NormalTok{) }\OperatorTok{{-}} \FloatTok{1}\OperatorTok{;}
\VariableTok{noise} \OperatorTok{=}\NormalTok{ (}\FloatTok{1}\OperatorTok{/}\VariableTok{sqrt}\NormalTok{(}\FloatTok{2}\OperatorTok{*}\VariableTok{EbN0}\NormalTok{)) }\OperatorTok{*} \VariableTok{randn}\NormalTok{(}\FloatTok{1}\OperatorTok{,} \VariableTok{N\_bits}\NormalTok{)}\OperatorTok{;}
\VariableTok{received} \OperatorTok{=} \VariableTok{bits} \OperatorTok{+} \VariableTok{noise}\OperatorTok{;}
\VariableTok{decisions} \OperatorTok{=} \VariableTok{sign}\NormalTok{(}\VariableTok{received}\NormalTok{)}\OperatorTok{;}
\VariableTok{errors} \OperatorTok{=} \VariableTok{sum}\NormalTok{(}\VariableTok{decisions} \OperatorTok{\textasciitilde{}=} \VariableTok{bits}\NormalTok{)}\OperatorTok{;}

\VariableTok{BER\_hat} \OperatorTok{=} \VariableTok{errors} \OperatorTok{/} \VariableTok{N\_bits}\OperatorTok{;}
\VariableTok{BER\_theory} \OperatorTok{=} \VariableTok{qfunc}\NormalTok{(}\VariableTok{sqrt}\NormalTok{(}\FloatTok{2}\OperatorTok{*}\VariableTok{EbN0}\NormalTok{))}\OperatorTok{;}

\CommentTok{\% 95\% Confidence Interval}
\VariableTok{margin} \OperatorTok{=} \FloatTok{1.96} \OperatorTok{*} \VariableTok{sqrt}\NormalTok{(}\VariableTok{BER\_hat}\OperatorTok{*}\NormalTok{(}\FloatTok{1}\OperatorTok{{-}}\VariableTok{BER\_hat}\NormalTok{)}\OperatorTok{/}\VariableTok{N\_bits}\NormalTok{)}\OperatorTok{;}
\VariableTok{fprintf}\NormalTok{(}\SpecialStringTok{\textquotesingle{}BER estimate: \%.4e\textbackslash{}n\textquotesingle{}}\OperatorTok{,} \VariableTok{BER\_hat}\NormalTok{)}\OperatorTok{;}
\VariableTok{fprintf}\NormalTok{(}\SpecialStringTok{\textquotesingle{}BER theory:   \%.4e\textbackslash{}n\textquotesingle{}}\OperatorTok{,} \VariableTok{BER\_theory}\NormalTok{)}\OperatorTok{;}
\VariableTok{fprintf}\NormalTok{(}\SpecialStringTok{\textquotesingle{}95\%\% CI: [\%.4e, \%.4e]\textbackslash{}n\textquotesingle{}}\OperatorTok{,} \VariableTok{BER\_hat}\OperatorTok{{-}}\VariableTok{margin}\OperatorTok{,} \VariableTok{BER\_hat}\OperatorTok{+}\VariableTok{margin}\NormalTok{)}\OperatorTok{;}
\end{Highlighting}
\end{Shaded}

\PSsubsection{16.1.2 AWGN Noise Modeling}

Additive White Gaussian Noise is the fundamental noise model:

\begin{itemize}
\tightlist
\item
  \textbf{Distribution:} $N\left(0,\, \frac{N_{0}}{2}\right)$ per dimension
\item
  \textbf{Why Gaussian?} CLT — thermal noise is sum of many electron contributions
\item
  \textbf{Characterization:} Fully described by power spectral density $N_{0}$ or variance $\sigma^{2} = N_{0}B$ (B = bandwidth)
\end{itemize}

\PSsubsection{16.1.3 Rayleigh Fading Channel}

\textbf{Model:} Channel coefficient $h = h_{I} + j \cdot h_{Q}$ where $h_{I},\, h_{Q} \sim N(0,\, \sigma^{2})$ independent.

\begin{itemize}
\tightlist
\item
  Envelope $\lvert h\rvert \sim \mathrm{Rayleigh}(\sigma)$
\item
  Power $\lvert h\rvert^{2} \sim \mathrm{Exponential}\left(\frac{1}{2\sigma^{2}}\right)$
\item
  \textbf{Outage probability:} $P(\lvert h\rvert^{2} < \gamma_{\mathrm{th}}) = 1 - \exp \left(-\frac{\gamma_{\mathrm{th}}}{2\sigma^{2}}\right)$
\end{itemize}

\tcbset{PScodemode/.style={unbreakable}}

\begin{Shaded}
\begin{Highlighting}[]
\CommentTok{\%\% Rayleigh Fading: Outage Probability}
\VariableTok{sigma} \OperatorTok{=} \FloatTok{1}\OperatorTok{;}
\VariableTok{SNR\_avg\_dB} \OperatorTok{=} \FloatTok{10}\OperatorTok{;}    \CommentTok{\% Average SNR in dB}
\VariableTok{SNR\_avg} \OperatorTok{=} \FloatTok{10}\OperatorTok{\^{}}\NormalTok{(}\VariableTok{SNR\_avg\_dB}\OperatorTok{/}\FloatTok{10}\NormalTok{)}\OperatorTok{;}
\VariableTok{SNR\_threshold\_dB} \OperatorTok{=} \FloatTok{5}\OperatorTok{;}  \CommentTok{\% Minimum required SNR}
\VariableTok{gamma\_th} \OperatorTok{=} \FloatTok{10}\OperatorTok{\^{}}\NormalTok{(}\VariableTok{SNR\_threshold\_dB}\OperatorTok{/}\FloatTok{10}\NormalTok{)}\OperatorTok{;}

\CommentTok{\% Analytical outage probability}
\VariableTok{P\_outage} \OperatorTok{=} \FloatTok{1} \OperatorTok{{-}} \VariableTok{exp}\NormalTok{(}\OperatorTok{{-}}\VariableTok{gamma\_th} \OperatorTok{/} \VariableTok{SNR\_avg}\NormalTok{)}\OperatorTok{;}
\VariableTok{fprintf}\NormalTok{(}\SpecialStringTok{\textquotesingle{}Outage Probability (analytical): \%.4f\textbackslash{}n\textquotesingle{}}\OperatorTok{,} \VariableTok{P\_outage}\NormalTok{)}\OperatorTok{;}

\CommentTok{\% Simulation}
\VariableTok{N} \OperatorTok{=} \FloatTok{100000}\OperatorTok{;}
\VariableTok{h} \OperatorTok{=}\NormalTok{ (}\VariableTok{randn}\NormalTok{(}\FloatTok{1}\OperatorTok{,}\VariableTok{N}\NormalTok{) }\OperatorTok{+} \FloatTok{1j}\OperatorTok{*}\VariableTok{randn}\NormalTok{(}\FloatTok{1}\OperatorTok{,}\VariableTok{N}\NormalTok{)) }\OperatorTok{/} \VariableTok{sqrt}\NormalTok{(}\FloatTok{2}\NormalTok{)}\OperatorTok{;}  \CommentTok{\% Rayleigh fading}
\VariableTok{SNR\_inst} \OperatorTok{=} \VariableTok{SNR\_avg} \OperatorTok{*} \VariableTok{abs}\NormalTok{(}\VariableTok{h}\NormalTok{)}\OperatorTok{.\^{}}\FloatTok{2}\OperatorTok{;}
\VariableTok{P\_outage\_sim} \OperatorTok{=} \VariableTok{mean}\NormalTok{(}\VariableTok{SNR\_inst} \OperatorTok{\textless{}} \VariableTok{gamma\_th}\NormalTok{)}\OperatorTok{;}
\VariableTok{fprintf}\NormalTok{(}\SpecialStringTok{\textquotesingle{}Outage Probability (simulated): \%.4f\textbackslash{}n\textquotesingle{}}\OperatorTok{,} \VariableTok{P\_outage\_sim}\NormalTok{)}\OperatorTok{;}
\end{Highlighting}
\end{Shaded}

\PSsubsection{16.1.4 Rician Fading}

When a line-of-sight (LOS) component exists:

\begin{itemize}
\tightlist
\item
  $h = h_{\mathrm{LOS}} + h_{\mathrm{scatter}}$
\item
  $\lvert h\rvert \sim \mathrm{Rician}(\nu,\, \sigma)$ where $K = \frac{\nu^{2}}{2\sigma^{2}}$ is the K-factor
\end{itemize}

K-factor: $K = 0$ (Rayleigh) $\to K = \infty$ (no fading, AWGN)

\PSsubsection{16.1.5 SNR Statistics}

\textbf{Instantaneous SNR:} $\gamma = \lvert h\rvert^{2} \times \frac{E_{s}}{N_{0}}$

For Rayleigh: $\gamma \sim \mathrm{Exponential}(\bar{\gamma})$ where $\bar{\gamma}$ = average SNR

\textbf{Average BER over fading:}\PSdisplay{\bar{P}_e = \int_0^\infty P_e(\gamma) \cdot f_\gamma(\gamma) \, d\gamma}

For BPSK over Rayleigh: $\bar{P}_{e} = \frac{1}{2}\left(1 - \sqrt{\frac{\bar{\gamma}}{1+\bar{\gamma}}}\right)$

\PSsection{16.2}{Signal Processing}

\PSsubsection{16.2.1 Measurement Noise Characterization}

\textbf{Problem:} Characterize noise in measured signal $x(t) = s(t) + n(t)$.

\textbf{Approach:}

\begin{enumerate}
\def\labelenumi{\arabic{enumi}.}
\tightlist
\item
  Collect noise-only samples (signal absent)
\item
  Estimate noise statistics: $\hat{\mu}_{n},\, \hat{\sigma}^{2}_{n}$
\item
  Test for Gaussianity (histogram, QQ-plot)
\item
  Compute SNR: $\mathrm{SNR} = \frac{P_{\mathrm{signal}}}{\hat{\sigma}^{2}_{n}}$
\end{enumerate}

\tcbset{PScodemode/.style={unbreakable}}

\begin{Shaded}
\begin{Highlighting}[]
\CommentTok{\%\% Noise Characterization}
\VariableTok{fs} \OperatorTok{=} \FloatTok{1000}\OperatorTok{;}  \VariableTok{T} \OperatorTok{=} \FloatTok{10}\OperatorTok{;}  \CommentTok{\% 1kHz sample rate, 10 seconds}
\VariableTok{t} \OperatorTok{=} \FloatTok{0}\OperatorTok{:}\FloatTok{1}\OperatorTok{/}\VariableTok{fs}\OperatorTok{:}\VariableTok{T}\OperatorTok{{-}}\FloatTok{1}\OperatorTok{/}\VariableTok{fs}\OperatorTok{;}
\VariableTok{N\_samples} \OperatorTok{=} \VariableTok{length}\NormalTok{(}\VariableTok{t}\NormalTok{)}\OperatorTok{;}

\CommentTok{\% Simulate: known signal + noise}
\VariableTok{signal} \OperatorTok{=} \FloatTok{2}\OperatorTok{*}\VariableTok{sin}\NormalTok{(}\FloatTok{2}\OperatorTok{*}\VariableTok{pi}\OperatorTok{*}\FloatTok{5}\OperatorTok{*}\VariableTok{t}\NormalTok{)}\OperatorTok{;}  \CommentTok{\% 5 Hz sine wave}
\VariableTok{noise} \OperatorTok{=} \FloatTok{0.5}\OperatorTok{*}\VariableTok{randn}\NormalTok{(}\FloatTok{1}\OperatorTok{,} \VariableTok{N\_samples}\NormalTok{)}\OperatorTok{;}
\VariableTok{measured} \OperatorTok{=} \VariableTok{signal} \OperatorTok{+} \VariableTok{noise}\OperatorTok{;}

\CommentTok{\% Estimate noise from signal{-}free segments (or residuals)}
\CommentTok{\% Here: subtract known signal}
\VariableTok{noise\_estimate} \OperatorTok{=} \VariableTok{measured} \OperatorTok{{-}} \VariableTok{signal}\OperatorTok{;}

\VariableTok{fprintf}\NormalTok{(}\SpecialStringTok{\textquotesingle{}Noise Statistics:\textbackslash{}n\textquotesingle{}}\NormalTok{)}\OperatorTok{;}
\VariableTok{fprintf}\NormalTok{(}\SpecialStringTok{\textquotesingle{}  Mean: \%.4f (should be \textasciitilde{}0)\textbackslash{}n\textquotesingle{}}\OperatorTok{,} \VariableTok{mean}\NormalTok{(}\VariableTok{noise\_estimate}\NormalTok{))}\OperatorTok{;}
\VariableTok{fprintf}\NormalTok{(}\SpecialStringTok{\textquotesingle{}  Std:  \%.4f (true: 0.5)\textbackslash{}n\textquotesingle{}}\OperatorTok{,} \VariableTok{std}\NormalTok{(}\VariableTok{noise\_estimate}\NormalTok{))}\OperatorTok{;}
\VariableTok{fprintf}\NormalTok{(}\SpecialStringTok{\textquotesingle{}  SNR:  \%.1f dB\textbackslash{}n\textquotesingle{}}\OperatorTok{,} \FloatTok{10}\OperatorTok{*}\VariableTok{log10}\NormalTok{(}\VariableTok{var}\NormalTok{(}\VariableTok{signal}\NormalTok{)}\OperatorTok{/}\VariableTok{var}\NormalTok{(}\VariableTok{noise\_estimate}\NormalTok{)))}\OperatorTok{;}

\CommentTok{\% Gaussianity test}
\VariableTok{figure}\OperatorTok{;}
\VariableTok{subplot}\NormalTok{(}\FloatTok{1}\OperatorTok{,}\FloatTok{2}\OperatorTok{,}\FloatTok{1}\NormalTok{)}\OperatorTok{;} \VariableTok{histogram}\NormalTok{(}\VariableTok{noise\_estimate}\OperatorTok{,} \FloatTok{50}\OperatorTok{,} \SpecialStringTok{\textquotesingle{}Normalization\textquotesingle{}}\OperatorTok{,} \SpecialStringTok{\textquotesingle{}pdf\textquotesingle{}}\NormalTok{)}\OperatorTok{;}
\VariableTok{hold} \VariableTok{on}\OperatorTok{;} \VariableTok{x} \OperatorTok{=} \VariableTok{linspace}\NormalTok{(}\OperatorTok{{-}}\FloatTok{2}\OperatorTok{,}\FloatTok{2}\OperatorTok{,}\FloatTok{200}\NormalTok{)}\OperatorTok{;} \VariableTok{plot}\NormalTok{(}\VariableTok{x}\OperatorTok{,} \VariableTok{normpdf}\NormalTok{(}\VariableTok{x}\OperatorTok{,}\FloatTok{0}\OperatorTok{,}\FloatTok{0.5}\NormalTok{)}\OperatorTok{,}\SpecialStringTok{\textquotesingle{}r\textquotesingle{}}\OperatorTok{,}\SpecialStringTok{\textquotesingle{}LineWidth\textquotesingle{}}\OperatorTok{,}\FloatTok{2}\NormalTok{)}\OperatorTok{;}
\VariableTok{title}\NormalTok{(}\SpecialStringTok{\textquotesingle{}Noise Histogram vs Gaussian\textquotesingle{}}\NormalTok{)}\OperatorTok{;}
\VariableTok{subplot}\NormalTok{(}\FloatTok{1}\OperatorTok{,}\FloatTok{2}\OperatorTok{,}\FloatTok{2}\NormalTok{)}\OperatorTok{;} \VariableTok{qqplot}\NormalTok{(}\VariableTok{noise\_estimate}\NormalTok{)}\OperatorTok{;} \VariableTok{title}\NormalTok{(}\SpecialStringTok{\textquotesingle{}QQ Plot\textquotesingle{}}\NormalTok{)}\OperatorTok{;}
\end{Highlighting}
\end{Shaded}

\PSsubsection{16.2.2 Correlation-Based Detection}

\textbf{Problem:} Detect a known signal in noise.

\textbf{Method:} Correlate received signal with known template. Under $H_{0}$ (noise only), the correlator output has known distribution \ensuremath{\to} set threshold.

\tcbset{PScodemode/.style={unbreakable}}

\begin{Shaded}
\begin{Highlighting}[]
\CommentTok{\%\% Correlation Detector}
\VariableTok{N} \OperatorTok{=} \FloatTok{100}\OperatorTok{;}  \CommentTok{\% Signal length}
\VariableTok{s} \OperatorTok{=} \VariableTok{ones}\NormalTok{(}\FloatTok{1}\OperatorTok{,}\VariableTok{N}\NormalTok{)}\OperatorTok{/}\VariableTok{sqrt}\NormalTok{(}\VariableTok{N}\NormalTok{)}\OperatorTok{;}  \CommentTok{\% Known signal template (normalized)}
\VariableTok{sigma\_n} \OperatorTok{=} \FloatTok{1}\OperatorTok{;}

\CommentTok{\% H0: noise only; H1: signal + noise}
\VariableTok{N\_trials} \OperatorTok{=} \FloatTok{10000}\OperatorTok{;}
\VariableTok{y\_H0} \OperatorTok{=} \VariableTok{zeros}\NormalTok{(}\FloatTok{1}\OperatorTok{,} \VariableTok{N\_trials}\NormalTok{)}\OperatorTok{;}  \VariableTok{y\_H1} \OperatorTok{=} \VariableTok{zeros}\NormalTok{(}\FloatTok{1}\OperatorTok{,} \VariableTok{N\_trials}\NormalTok{)}\OperatorTok{;}
\VariableTok{A} \OperatorTok{=} \FloatTok{2}\OperatorTok{;}  \CommentTok{\% Signal amplitude}

\KeywordTok{for} \VariableTok{i} \OperatorTok{=} \FloatTok{1}\OperatorTok{:}\VariableTok{N\_trials}
    \VariableTok{n} \OperatorTok{=} \VariableTok{sigma\_n} \OperatorTok{*} \VariableTok{randn}\NormalTok{(}\FloatTok{1}\OperatorTok{,} \VariableTok{N}\NormalTok{)}\OperatorTok{;}
    \VariableTok{y\_H0}\NormalTok{(}\VariableTok{i}\NormalTok{) }\OperatorTok{=} \VariableTok{s} \OperatorTok{*} \VariableTok{n}\OperatorTok{\textquotesingle{};}              \CommentTok{\% Correlator output under H0}
    \VariableTok{y\_H1}\NormalTok{(}\VariableTok{i}\NormalTok{) }\OperatorTok{=} \VariableTok{s} \OperatorTok{*}\NormalTok{ (}\VariableTok{A}\OperatorTok{*}\VariableTok{s} \OperatorTok{+} \VariableTok{n}\NormalTok{)}\OperatorTok{\textquotesingle{};}     \CommentTok{\% Correlator output under H1}
\KeywordTok{end}

\CommentTok{\% Set threshold for P\_FA = 0.01}
\VariableTok{threshold} \OperatorTok{=} \VariableTok{sigma\_n}\OperatorTok{/}\VariableTok{sqrt}\NormalTok{(}\VariableTok{N}\NormalTok{) }\OperatorTok{*} \VariableTok{norminv}\NormalTok{(}\FloatTok{0.99}\NormalTok{)}\OperatorTok{;}
\VariableTok{P\_FA} \OperatorTok{=} \VariableTok{mean}\NormalTok{(}\VariableTok{y\_H0} \OperatorTok{\textgreater{}} \VariableTok{threshold}\NormalTok{)}\OperatorTok{;}
\VariableTok{P\_D} \OperatorTok{=} \VariableTok{mean}\NormalTok{(}\VariableTok{y\_H1} \OperatorTok{\textgreater{}} \VariableTok{threshold}\NormalTok{)}\OperatorTok{;}
\VariableTok{fprintf}\NormalTok{(}\SpecialStringTok{\textquotesingle{}P\_FA = \%.4f (target: 0.01)\textbackslash{}n\textquotesingle{}}\OperatorTok{,} \VariableTok{P\_FA}\NormalTok{)}\OperatorTok{;}
\VariableTok{fprintf}\NormalTok{(}\SpecialStringTok{\textquotesingle{}P\_D = \%.4f\textbackslash{}n\textquotesingle{}}\OperatorTok{,} \VariableTok{P\_D}\NormalTok{)}\OperatorTok{;}
\end{Highlighting}
\end{Shaded}

\PSsection{16.3}{Electronics}

\PSsubsection{16.3.1 Component Tolerance Modeling}

\textbf{Problem:} Resistors have nominal value $R_{0}$ with manufacturing tolerance.

\textbf{Model:} $R \sim N(R_{0},\, \sigma_{R}^{2})$ where $\sigma_{R}$ depends on tolerance class.

For $\pm 5\%$ tolerance (99.7\% within bounds): $3\sigma = 0.05R_{0} \to \sigma = 0.0167R_{0}$

\PSsubsection{16.3.2 Monte Carlo Circuit Analysis}

\textbf{Problem:} A voltage divider uses $R_{1}$ and $R_{2}$. What is the distribution of output voltage?

$\displaystyle V_{\mathrm{out}} = V_{\mathrm{in}} \times \frac{R_{2}}{R_{1} + R_{2}}$

With $R_{1},\, R_{2}$ random: $V_{\mathrm{out}}$ is a random variable!

\tcbset{PScodemode/.style={unbreakable}}

\begin{Shaded}
\begin{Highlighting}[]
\CommentTok{\%\% Monte Carlo: Voltage Divider with Tolerant Components}
\VariableTok{V\_in} \OperatorTok{=} \FloatTok{5}\OperatorTok{;}         \CommentTok{\% Input voltage}
\VariableTok{R1\_nom} \OperatorTok{=} \FloatTok{10e3}\OperatorTok{;}    \VariableTok{R2\_nom} \OperatorTok{=} \FloatTok{10e3}\OperatorTok{;}   \CommentTok{\% 10kΩ each}
\VariableTok{tol} \OperatorTok{=} \FloatTok{0.05}\OperatorTok{;}       \CommentTok{\% 5\% tolerance}
\VariableTok{sigma\_R} \OperatorTok{=} \VariableTok{tol} \OperatorTok{*} \VariableTok{R1\_nom} \OperatorTok{/} \FloatTok{3}\OperatorTok{;}  \CommentTok{\% 3{-}sigma = 5\%}

\VariableTok{N} \OperatorTok{=} \FloatTok{100000}\OperatorTok{;}
\VariableTok{R1} \OperatorTok{=} \VariableTok{R1\_nom} \OperatorTok{+} \VariableTok{sigma\_R}\OperatorTok{*}\VariableTok{randn}\NormalTok{(}\FloatTok{1}\OperatorTok{,} \VariableTok{N}\NormalTok{)}\OperatorTok{;}
\VariableTok{R2} \OperatorTok{=} \VariableTok{R2\_nom} \OperatorTok{+} \VariableTok{sigma\_R}\OperatorTok{*}\VariableTok{randn}\NormalTok{(}\FloatTok{1}\OperatorTok{,} \VariableTok{N}\NormalTok{)}\OperatorTok{;}
\VariableTok{V\_out} \OperatorTok{=} \VariableTok{V\_in} \OperatorTok{*} \VariableTok{R2} \OperatorTok{./}\NormalTok{ (}\VariableTok{R1} \OperatorTok{+} \VariableTok{R2}\NormalTok{)}\OperatorTok{;}

\VariableTok{fprintf}\NormalTok{(}\SpecialStringTok{\textquotesingle{}Nominal V\_out = \%.3f V\textbackslash{}n\textquotesingle{}}\OperatorTok{,} \VariableTok{V\_in} \OperatorTok{*} \VariableTok{R2\_nom}\OperatorTok{/}\NormalTok{(}\VariableTok{R1\_nom}\OperatorTok{+}\VariableTok{R2\_nom}\NormalTok{))}\OperatorTok{;}
\VariableTok{fprintf}\NormalTok{(}\SpecialStringTok{\textquotesingle{}Mean V\_out = \%.4f V\textbackslash{}n\textquotesingle{}}\OperatorTok{,} \VariableTok{mean}\NormalTok{(}\VariableTok{V\_out}\NormalTok{))}\OperatorTok{;}
\VariableTok{fprintf}\NormalTok{(}\SpecialStringTok{\textquotesingle{}Std V\_out = \%.4f V\textbackslash{}n\textquotesingle{}}\OperatorTok{,} \VariableTok{std}\NormalTok{(}\VariableTok{V\_out}\NormalTok{))}\OperatorTok{;}
\VariableTok{fprintf}\NormalTok{(}\SpecialStringTok{\textquotesingle{}99.7\%\% range: [\%.3f, \%.3f] V\textbackslash{}n\textquotesingle{}}\OperatorTok{,} \OperatorTok{...}
    \VariableTok{mean}\NormalTok{(}\VariableTok{V\_out}\NormalTok{)}\OperatorTok{{-}}\FloatTok{3}\OperatorTok{*}\VariableTok{std}\NormalTok{(}\VariableTok{V\_out}\NormalTok{)}\OperatorTok{,} \VariableTok{mean}\NormalTok{(}\VariableTok{V\_out}\NormalTok{)}\OperatorTok{+}\FloatTok{3}\OperatorTok{*}\VariableTok{std}\NormalTok{(}\VariableTok{V\_out}\NormalTok{))}\OperatorTok{;}

\CommentTok{\% Yield: fraction within ±2\% of nominal}
\VariableTok{V\_nom} \OperatorTok{=} \FloatTok{2.5}\OperatorTok{;}
\VariableTok{yield} \OperatorTok{=} \VariableTok{mean}\NormalTok{(}\VariableTok{abs}\NormalTok{(}\VariableTok{V\_out} \OperatorTok{{-}} \VariableTok{V\_nom}\NormalTok{) }\OperatorTok{\textless{}} \FloatTok{0.02}\OperatorTok{*}\VariableTok{V\_nom}\NormalTok{)}\OperatorTok{;}
\VariableTok{fprintf}\NormalTok{(}\SpecialStringTok{\textquotesingle{}Yield (±2\%\%): \%.2f\%\%\textbackslash{}n\textquotesingle{}}\OperatorTok{,} \FloatTok{100}\OperatorTok{*}\VariableTok{yield}\NormalTok{)}\OperatorTok{;}
\end{Highlighting}
\end{Shaded}

\PSsubsection{16.3.3 Measurement Uncertainty Propagation}

If $Z = f(X,\, Y)$ and $X,\, Y$ have small uncertainties:

\PSdisplay{\sigma_Z^2 \approx \left(\frac{\partial f}{\partial X}\right)^2 \sigma_X^2 + \left(\frac{\partial f}{\partial Y}\right)^2 \sigma_Y^2}

\textbf{Example:} Power $P = \frac{V^{2}}{R}$. $\sigma_{P}^{2} \approx \frac{2\,\mathrm{V}}{R}^{2}\sigma_{V}^{2} + \frac{V^{2}}{R^{2}}^{2}\sigma_{R}^{2}$

\PSsection{16.4}{Power Systems}

\PSsubsection{16.4.1 Load Uncertainty}

\textbf{Problem:} Peak load is uncertain. Model: $L \sim N(\mu_{L},\, \sigma_{L}^{2})$

Design question: What capacity $C$ ensures $P(\text{load} > \text{capacity}) < 0.01$?

$C = \mu_{L} + 2.326\sigma_{L}$ (for 1\% exceedance probability)

\PSsubsection{16.4.2 System Reliability}

\textbf{Series system:} All components must work.

$R_{\mathrm{sys}} = R_{1} \times R_{2} \times \ldots \times R_{n}$ (for independent components)

\textbf{Parallel system:} At least one must work.

$R_{\mathrm{sys}} = 1 - (1-R_{1})(1-R_{2})\ldots (1-R_{n})$

\tcbset{PScodemode/.style={unbreakable}}

\begin{Shaded}
\begin{Highlighting}[]
\CommentTok{\%\% Reliability Analysis: Series{-}Parallel System}
\CommentTok{\% System: Two subsystems in series, each with 3 parallel redundant units}
\VariableTok{R\_unit} \OperatorTok{=} \FloatTok{0.95}\OperatorTok{;}  \CommentTok{\% Individual component reliability}
\VariableTok{n\_parallel} \OperatorTok{=} \FloatTok{3}\OperatorTok{;}
\VariableTok{n\_series} \OperatorTok{=} \FloatTok{2}\OperatorTok{;}

\CommentTok{\% Each parallel subsystem}
\VariableTok{R\_subsystem} \OperatorTok{=} \FloatTok{1} \OperatorTok{{-}}\NormalTok{ (}\FloatTok{1} \OperatorTok{{-}} \VariableTok{R\_unit}\NormalTok{)}\OperatorTok{\^{}}\VariableTok{n\_parallel}\OperatorTok{;}
\CommentTok{\% Series combination}
\VariableTok{R\_system} \OperatorTok{=} \VariableTok{R\_subsystem}\OperatorTok{\^{}}\VariableTok{n\_series}\OperatorTok{;}

\VariableTok{fprintf}\NormalTok{(}\SpecialStringTok{\textquotesingle{}Unit reliability: \%.4f\textbackslash{}n\textquotesingle{}}\OperatorTok{,} \VariableTok{R\_unit}\NormalTok{)}\OperatorTok{;}
\VariableTok{fprintf}\NormalTok{(}\SpecialStringTok{\textquotesingle{}Subsystem (3 parallel): \%.6f\textbackslash{}n\textquotesingle{}}\OperatorTok{,} \VariableTok{R\_subsystem}\NormalTok{)}\OperatorTok{;}
\VariableTok{fprintf}\NormalTok{(}\SpecialStringTok{\textquotesingle{}System (2 in series): \%.6f\textbackslash{}n\textquotesingle{}}\OperatorTok{,} \VariableTok{R\_system}\NormalTok{)}\OperatorTok{;}

\CommentTok{\% Monte Carlo verification}
\VariableTok{N} \OperatorTok{=} \FloatTok{100000}\OperatorTok{;}
\VariableTok{sys\_works} \OperatorTok{=} \VariableTok{true}\NormalTok{(}\FloatTok{1}\OperatorTok{,} \VariableTok{N}\NormalTok{)}\OperatorTok{;}
\KeywordTok{for} \VariableTok{s} \OperatorTok{=} \FloatTok{1}\OperatorTok{:}\VariableTok{n\_series}
    \VariableTok{sub\_works} \OperatorTok{=} \VariableTok{false}\NormalTok{(}\FloatTok{1}\OperatorTok{,} \VariableTok{N}\NormalTok{)}\OperatorTok{;}
    \KeywordTok{for} \VariableTok{p} \OperatorTok{=} \FloatTok{1}\OperatorTok{:}\VariableTok{n\_parallel}
        \VariableTok{unit\_works} \OperatorTok{=} \VariableTok{rand}\NormalTok{(}\FloatTok{1}\OperatorTok{,} \VariableTok{N}\NormalTok{) }\OperatorTok{\textless{}} \VariableTok{R\_unit}\OperatorTok{;}
        \VariableTok{sub\_works} \OperatorTok{=} \VariableTok{sub\_works} \OperatorTok{|} \VariableTok{unit\_works}\OperatorTok{;}
    \KeywordTok{end}
    \VariableTok{sys\_works} \OperatorTok{=} \VariableTok{sys\_works} \OperatorTok{\&} \VariableTok{sub\_works}\OperatorTok{;}
\KeywordTok{end}
\VariableTok{fprintf}\NormalTok{(}\SpecialStringTok{\textquotesingle{}System reliability (sim): \%.6f\textbackslash{}n\textquotesingle{}}\OperatorTok{,} \VariableTok{mean}\NormalTok{(}\VariableTok{sys\_works}\NormalTok{))}\OperatorTok{;}
\end{Highlighting}
\end{Shaded}

\PSsection{16.5}{Control Systems}

\PSsubsection{16.5.1 Sensor Noise Effects}

\textbf{Problem:} A control system uses noisy sensor readings.

Measurement: $y(k) = x(k) + n(k)$ where $n \sim N(0,\, \sigma^{2}_{n})$

Controller acts on $y(k)$ \ensuremath{\to} tracking error includes noise contribution.

\textbf{Solution approaches:}

\begin{itemize}
\tightlist
\item
  Averaging (low-pass filter): reduces noise by $\sqrt{n}$
\item
  Kalman filter: optimal fusion of model prediction and measurement
\end{itemize}

\tcbset{PScodemode/.style={unbreakable}}

\begin{Shaded}
\begin{Highlighting}[]
\CommentTok{\%\% Noisy Sensor in Control Loop: Effect of Averaging}
\VariableTok{true\_value} \OperatorTok{=} \FloatTok{10}\OperatorTok{;}    \CommentTok{\% Constant being measured}
\VariableTok{sigma\_n} \OperatorTok{=} \FloatTok{0.5}\OperatorTok{;}      \CommentTok{\% Sensor noise std dev}
\VariableTok{N\_measurements} \OperatorTok{=} \FloatTok{1000}\OperatorTok{;}

\CommentTok{\% Raw measurements}
\VariableTok{measurements} \OperatorTok{=} \VariableTok{true\_value} \OperatorTok{+} \VariableTok{sigma\_n}\OperatorTok{*}\VariableTok{randn}\NormalTok{(}\FloatTok{1}\OperatorTok{,} \VariableTok{N\_measurements}\NormalTok{)}\OperatorTok{;}

\CommentTok{\% Moving average filter (window size M)}
\VariableTok{M\_values} \OperatorTok{=}\NormalTok{ [}\FloatTok{1}\OperatorTok{,} \FloatTok{5}\OperatorTok{,} \FloatTok{10}\OperatorTok{,} \FloatTok{20}\OperatorTok{,} \FloatTok{50}\NormalTok{]}\OperatorTok{;}
\VariableTok{figure}\OperatorTok{;} \VariableTok{hold} \VariableTok{on}\OperatorTok{;}
\KeywordTok{for} \VariableTok{M} \OperatorTok{=} \VariableTok{M\_values}
    \VariableTok{filtered} \OperatorTok{=} \VariableTok{movmean}\NormalTok{(}\VariableTok{measurements}\OperatorTok{,} \VariableTok{M}\NormalTok{)}\OperatorTok{;}
    \VariableTok{plot}\NormalTok{(}\VariableTok{filtered}\NormalTok{(}\FloatTok{100}\OperatorTok{:}\FloatTok{500}\NormalTok{))}\OperatorTok{;}
\KeywordTok{end}
\VariableTok{yline}\NormalTok{(}\VariableTok{true\_value}\OperatorTok{,} \SpecialStringTok{\textquotesingle{}k{-}{-}\textquotesingle{}}\OperatorTok{,} \SpecialStringTok{\textquotesingle{}LineWidth\textquotesingle{}}\OperatorTok{,} \FloatTok{2}\NormalTok{)}\OperatorTok{;}
\VariableTok{legend}\NormalTok{(}\VariableTok{arrayfun}\NormalTok{(}\OperatorTok{@}\NormalTok{(}\VariableTok{m}\NormalTok{) }\VariableTok{sprintf}\NormalTok{(}\SpecialStringTok{\textquotesingle{}M=\%d, σ=\%.3f\textquotesingle{}}\OperatorTok{,} \VariableTok{m}\OperatorTok{,} \VariableTok{sigma\_n}\OperatorTok{/}\VariableTok{sqrt}\NormalTok{(}\VariableTok{m}\NormalTok{))}\OperatorTok{,} \OperatorTok{...}
    \VariableTok{M\_values}\OperatorTok{,} \SpecialStringTok{\textquotesingle{}UniformOutput\textquotesingle{}}\OperatorTok{,} \VariableTok{false}\NormalTok{))}\OperatorTok{;}
\VariableTok{xlabel}\NormalTok{(}\SpecialStringTok{\textquotesingle{}Sample\textquotesingle{}}\NormalTok{)}\OperatorTok{;} \VariableTok{ylabel}\NormalTok{(}\SpecialStringTok{\textquotesingle{}Estimate\textquotesingle{}}\NormalTok{)}\OperatorTok{;}
\VariableTok{title}\NormalTok{(}\SpecialStringTok{\textquotesingle{}Noise Reduction by Averaging\textquotesingle{}}\NormalTok{)}\OperatorTok{;}
\end{Highlighting}
\end{Shaded}

\PSsubsection{16.5.2 Parameter Uncertainty}

System model: $G(s) = \frac{K}{\tau s + 1}$ where $K$ and $\tau$ are uncertain.

Monte Carlo approach: sample $K$ and $\tau$ from their distributions, simulate system response many times, compute statistics of output.

\PSsection{16.6}{Machine Learning / AI Connections}

\PSsubsection{16.6.1 Sampling and Data Collection}

\begin{itemize}
\tightlist
\item
  Training data = sample from population
\item
  Test data = independent sample for validation
\item
  i.i.d. assumption underlies most ML theory
\end{itemize}

\PSsubsection{16.6.2 Estimation as Foundation of ML}

\begin{itemize}
\tightlist
\item
  Linear regression = MLE under Gaussian noise assumption
\item
  Neural network training = MLE (or MAP) for model parameters
\item
  Loss function minimization = maximum likelihood principle
\end{itemize}

\PSsubsection{16.6.3 Probability Distributions in ML}

\begin{itemize}
\tightlist
\item
  Generative models: learn $P(X)$ directly
\item
  Gaussian Mixture Models: weighted sum of Gaussians
\item
  Variational Autoencoders: latent variables with Gaussian priors
\end{itemize}

\PSsubsection{16.6.4 Bias-Variance Tradeoff}

\PSdisplay{\text{Test Error} = \text{Bias}^2 + \text{Variance} + \text{Irreducible Noise}}

This is exactly the MSE decomposition from estimation theory (Module 13)!

\begin{itemize}
\tightlist
\item
  High bias: underfitting (too simple model)
\item
  High variance: overfitting (too complex model)
\item
  Optimal: balance between the two
\end{itemize}

\PSsubsection{16.6.5 Confidence and Uncertainty in Predictions}

\begin{itemize}
\tightlist
\item
  Confidence intervals for model predictions
\item
  Bayesian neural networks: posterior over weights \ensuremath{\to} prediction uncertainty
\item
  Calibration: predicted probabilities should match observed frequencies
\end{itemize}

\PSsection{16.7}{Comprehensive MATLAB Project: Communication System Analysis}

\tcbset{PScodemode/.style={breakable}}

\begin{Shaded}
\begin{Highlighting}[]
\CommentTok{\%\% Complete Communication System Statistical Analysis}
\CommentTok{\% Simulate BPSK over Rayleigh fading with AWGN}
\CommentTok{\% Apply estimation, CI, and hypothesis testing}

\CommentTok{\% System parameters}
\VariableTok{N\_bits} \OperatorTok{=} \FloatTok{1e5}\OperatorTok{;}
\VariableTok{EbN0\_dB} \OperatorTok{=} \FloatTok{10}\OperatorTok{;}
\VariableTok{EbN0} \OperatorTok{=} \FloatTok{10}\OperatorTok{\^{}}\NormalTok{(}\VariableTok{EbN0\_dB}\OperatorTok{/}\FloatTok{10}\NormalTok{)}\OperatorTok{;}

\CommentTok{\% Generate channel and signal}
\VariableTok{bits} \OperatorTok{=} \VariableTok{randi}\NormalTok{([}\FloatTok{0}\OperatorTok{,}\FloatTok{1}\NormalTok{]}\OperatorTok{,} \FloatTok{1}\OperatorTok{,} \VariableTok{N\_bits}\NormalTok{)}\OperatorTok{;}
\VariableTok{symbols} \OperatorTok{=} \FloatTok{2}\OperatorTok{*}\VariableTok{bits} \OperatorTok{{-}} \FloatTok{1}\OperatorTok{;}  \CommentTok{\% BPSK: ±1}

\CommentTok{\% Rayleigh fading channel}
\VariableTok{h} \OperatorTok{=}\NormalTok{ (}\VariableTok{randn}\NormalTok{(}\FloatTok{1}\OperatorTok{,}\VariableTok{N\_bits}\NormalTok{) }\OperatorTok{+} \FloatTok{1j}\OperatorTok{*}\VariableTok{randn}\NormalTok{(}\FloatTok{1}\OperatorTok{,}\VariableTok{N\_bits}\NormalTok{)) }\OperatorTok{/} \VariableTok{sqrt}\NormalTok{(}\FloatTok{2}\NormalTok{)}\OperatorTok{;}
\VariableTok{h\_mag} \OperatorTok{=} \VariableTok{abs}\NormalTok{(}\VariableTok{h}\NormalTok{)}\OperatorTok{;}

\CommentTok{\% AWGN noise}
\VariableTok{noise\_var} \OperatorTok{=} \FloatTok{1}\OperatorTok{/}\NormalTok{(}\FloatTok{2}\OperatorTok{*}\VariableTok{EbN0}\NormalTok{)}\OperatorTok{;}
\VariableTok{noise} \OperatorTok{=} \VariableTok{sqrt}\NormalTok{(}\VariableTok{noise\_var}\NormalTok{) }\OperatorTok{*} \VariableTok{randn}\NormalTok{(}\FloatTok{1}\OperatorTok{,} \VariableTok{N\_bits}\NormalTok{)}\OperatorTok{;}

\CommentTok{\% Received signal (with fading)}
\VariableTok{received} \OperatorTok{=} \VariableTok{h\_mag} \OperatorTok{.*} \VariableTok{symbols} \OperatorTok{+} \VariableTok{noise}\OperatorTok{;}

\CommentTok{\% Detection (coherent, known channel)}
\VariableTok{detected} \OperatorTok{=} \VariableTok{sign}\NormalTok{(}\VariableTok{real}\NormalTok{(}\VariableTok{received} \OperatorTok{./} \VariableTok{h\_mag}\NormalTok{))}\OperatorTok{;}  \CommentTok{\% Zero{-}forcing}
\VariableTok{errors} \OperatorTok{=} \VariableTok{sum}\NormalTok{(}\VariableTok{detected} \OperatorTok{\textasciitilde{}=} \VariableTok{symbols}\NormalTok{)}\OperatorTok{;}
\VariableTok{BER\_measured} \OperatorTok{=} \VariableTok{errors} \OperatorTok{/} \VariableTok{N\_bits}\OperatorTok{;}

\CommentTok{\% {-}{-}{-} Statistical Analysis {-}{-}{-}}
\VariableTok{fprintf}\NormalTok{(}\SpecialStringTok{\textquotesingle{}=== Communication System Statistical Analysis ===\textbackslash{}n\textquotesingle{}}\NormalTok{)}\OperatorTok{;}
\VariableTok{fprintf}\NormalTok{(}\SpecialStringTok{\textquotesingle{}Eb/N0 = \%d dB, N = \%d bits\textbackslash{}n\textbackslash{}n\textquotesingle{}}\OperatorTok{,} \VariableTok{EbN0\_dB}\OperatorTok{,} \VariableTok{N\_bits}\NormalTok{)}\OperatorTok{;}

\CommentTok{\% 1. BER Estimation}
\VariableTok{fprintf}\NormalTok{(}\SpecialStringTok{\textquotesingle{}1. BER ESTIMATION\textbackslash{}n\textquotesingle{}}\NormalTok{)}\OperatorTok{;}
\VariableTok{fprintf}\NormalTok{(}\SpecialStringTok{\textquotesingle{}   Measured BER: \%.4e (\%d errors)\textbackslash{}n\textquotesingle{}}\OperatorTok{,} \VariableTok{BER\_measured}\OperatorTok{,} \VariableTok{errors}\NormalTok{)}\OperatorTok{;}
\VariableTok{BER\_theory\_rayleigh} \OperatorTok{=} \FloatTok{0.5}\OperatorTok{*}\NormalTok{(}\FloatTok{1} \OperatorTok{{-}} \VariableTok{sqrt}\NormalTok{(}\VariableTok{EbN0}\OperatorTok{/}\NormalTok{(}\FloatTok{1}\OperatorTok{+}\VariableTok{EbN0}\NormalTok{)))}\OperatorTok{;}
\VariableTok{fprintf}\NormalTok{(}\SpecialStringTok{\textquotesingle{}   Theory (Rayleigh): \%.4e\textbackslash{}n\textquotesingle{}}\OperatorTok{,} \VariableTok{BER\_theory\_rayleigh}\NormalTok{)}\OperatorTok{;}

\CommentTok{\% 2. Confidence Interval}
\VariableTok{fprintf}\NormalTok{(}\SpecialStringTok{\textquotesingle{}\textbackslash{}n2. CONFIDENCE INTERVAL (95\%\%)\textbackslash{}n\textquotesingle{}}\NormalTok{)}\OperatorTok{;}
\VariableTok{margin} \OperatorTok{=} \FloatTok{1.96}\OperatorTok{*}\VariableTok{sqrt}\NormalTok{(}\VariableTok{BER\_measured}\OperatorTok{*}\NormalTok{(}\FloatTok{1}\OperatorTok{{-}}\VariableTok{BER\_measured}\NormalTok{)}\OperatorTok{/}\VariableTok{N\_bits}\NormalTok{)}\OperatorTok{;}
\VariableTok{fprintf}\NormalTok{(}\SpecialStringTok{\textquotesingle{}   CI: [\%.4e, \%.4e]\textbackslash{}n\textquotesingle{}}\OperatorTok{,} \VariableTok{BER\_measured}\OperatorTok{{-}}\VariableTok{margin}\OperatorTok{,} \VariableTok{BER\_measured}\OperatorTok{+}\VariableTok{margin}\NormalTok{)}\OperatorTok{;}
\VariableTok{fprintf}\NormalTok{(}\SpecialStringTok{\textquotesingle{}   Theory in CI? \%s\textbackslash{}n\textquotesingle{}}\OperatorTok{,} \OperatorTok{...}
    \VariableTok{string}\NormalTok{(}\VariableTok{BER\_theory\_rayleigh} \OperatorTok{\textgreater{}=} \VariableTok{BER\_measured}\OperatorTok{{-}}\VariableTok{margin} \OperatorTok{\&} \OperatorTok{...}
           \VariableTok{BER\_theory\_rayleigh} \OperatorTok{\textless{}=} \VariableTok{BER\_measured}\OperatorTok{+}\VariableTok{margin}\NormalTok{))}\OperatorTok{;}

\CommentTok{\% 3. Channel Statistics}
\VariableTok{fprintf}\NormalTok{(}\SpecialStringTok{\textquotesingle{}\textbackslash{}n3. CHANNEL STATISTICS\textbackslash{}n\textquotesingle{}}\NormalTok{)}\OperatorTok{;}
\VariableTok{fprintf}\NormalTok{(}\SpecialStringTok{\textquotesingle{}   Mean |h|: \%.3f (theory: \%.3f)\textbackslash{}n\textquotesingle{}}\OperatorTok{,} \VariableTok{mean}\NormalTok{(}\VariableTok{h\_mag}\NormalTok{)}\OperatorTok{,} \VariableTok{sqrt}\NormalTok{(}\VariableTok{pi}\OperatorTok{/}\FloatTok{4}\NormalTok{))}\OperatorTok{;}
\VariableTok{fprintf}\NormalTok{(}\SpecialStringTok{\textquotesingle{}   Var |h|²: \%.3f (theory: \%.3f)\textbackslash{}n\textquotesingle{}}\OperatorTok{,} \VariableTok{var}\NormalTok{(}\VariableTok{h\_mag}\OperatorTok{.\^{}}\FloatTok{2}\NormalTok{)}\OperatorTok{,} \FloatTok{1}\OperatorTok{{-}}\NormalTok{(}\VariableTok{pi}\OperatorTok{/}\FloatTok{4}\NormalTok{))}\OperatorTok{;}

\CommentTok{\% 4. Outage Probability}
\VariableTok{SNR\_threshold} \OperatorTok{=} \FloatTok{3}\OperatorTok{;}  \CommentTok{\% dB}
\VariableTok{gamma\_th} \OperatorTok{=} \FloatTok{10}\OperatorTok{\^{}}\NormalTok{(}\VariableTok{SNR\_threshold}\OperatorTok{/}\FloatTok{10}\NormalTok{)}\OperatorTok{;}
\VariableTok{P\_outage\_sim} \OperatorTok{=} \VariableTok{mean}\NormalTok{(}\VariableTok{EbN0}\OperatorTok{*}\VariableTok{h\_mag}\OperatorTok{.\^{}}\FloatTok{2} \OperatorTok{\textless{}} \VariableTok{gamma\_th}\NormalTok{)}\OperatorTok{;}
\VariableTok{P\_outage\_theory} \OperatorTok{=} \FloatTok{1} \OperatorTok{{-}} \VariableTok{exp}\NormalTok{(}\OperatorTok{{-}}\VariableTok{gamma\_th}\OperatorTok{/}\VariableTok{EbN0}\NormalTok{)}\OperatorTok{;}
\VariableTok{fprintf}\NormalTok{(}\SpecialStringTok{\textquotesingle{}\textbackslash{}n4. OUTAGE PROBABILITY (threshold = \%d dB)\textbackslash{}n\textquotesingle{}}\OperatorTok{,} \VariableTok{SNR\_threshold}\NormalTok{)}\OperatorTok{;}
\VariableTok{fprintf}\NormalTok{(}\SpecialStringTok{\textquotesingle{}   Simulated: \%.4f\textbackslash{}n\textquotesingle{}}\OperatorTok{,} \VariableTok{P\_outage\_sim}\NormalTok{)}\OperatorTok{;}
\VariableTok{fprintf}\NormalTok{(}\SpecialStringTok{\textquotesingle{}   Theory:    \%.4f\textbackslash{}n\textquotesingle{}}\OperatorTok{,} \VariableTok{P\_outage\_theory}\NormalTok{)}\OperatorTok{;}
\end{Highlighting}
\end{Shaded}

\PSsection{16.8}{Key Integration Points}

{\def\LTcaptype{none} % do not increment counter
\begin{longtable}[c]{@{}cc@{}}
\toprule\noalign{}
\rowcolor{PSnavy}\PSth{Module} & \PSth{Application in EE} \\
\midrule\noalign{}
\endhead
\bottomrule\noalign{}
\endlastfoot
1. Probability & Event probabilities, fault diagnosis \\
2. Random Variables & Modeling uncertain quantities \\
3. Distributions & Choosing noise/fading/failure models \\
4. Moments & Signal power, noise power, SNR \\
5. Transformations & dB conversion, power from voltage \\
$6-7$. Joint/Functions & Signal+noise, system combinations \\
8. Conditional & Performance under different states \\
9. Correlation & MIMO channels, sensor arrays \\
10. Multiple RVs & Vector processing, beamforming \\
11. CLT & Why noise is Gaussian, averaging \\
12. Statistics & Estimating system parameters \\
13. Estimation & MLE for channel/noise parameters \\
14. Confidence & Quantifying measurement uncertainty \\
15. Hypothesis & Quality control, performance comparison \\
\end{longtable}
}

\emph{End of course modules. See {Supplementary Materials} for decision guides, MATLAB reference, and formula sheet.}
